\documentclass[a4paper,11pt]{article}
\usepackage[utf8]{inputenc}
\usepackage[T1]{fontenc}
\usepackage{lmodern}
\usepackage[french]{babel}
\usepackage{geometry}
\geometry{top=2.5cm, bottom=2.5cm, left=2.2cm, right=2.2cm}
\usepackage{xcolor}
\usepackage{titlesec}
\usepackage{enumitem}
\usepackage{booktabs}
\usepackage{tabularx}
\usepackage{fancyhdr}
\usepackage{hyperref}

% Définition des couleurs
\definecolor{enstablu}{RGB}{0, 82, 156} % Bleu institutionnel
\definecolor{darkgray}{RGB}{60, 60, 60}
\definecolor{lightgray}{RGB}{245, 245, 245}

% Configuration de hyperref
\hypersetup{
    colorlinks=true,
    linkcolor=enstablu,
    filecolor=enstablu,      
    urlcolor=enstablu,
    pdftitle={Cahier des charges - Système informatique des Foy'z/Bar},
    pdfauthor={Paul Debise}
}

% Configuration de la mise en page (fancyhdr)
\pagestyle{fancy}
\setlength{\headheight}{15pt}
\fancyhf{}
\fancyhead[L]{\textcolor{darkgray}{\small Cahier des charges | Système informatique des Foy'z/Bar}}
\fancyhead[R]{\textcolor{darkgray}{\small ENSTA Brest $\&$ Paris}}
\fancyfoot[C]{\thepage}
\renewcommand{\headrulewidth}{0.4pt}
\renewcommand{\footrulewidth}{0pt}

% Style des titres
\titleformat{\section}{\color{enstablu}\normalfont\Large\bfseries}{\thesection}{1em}{}[{\titlerule[0.8pt]}]
\titleformat{\subsection}{\color{enstablu}\normalfont\large\bfseries}{\thesubsection}{1em}{}
\titleformat{\subsubsection}{\color{darkgray}\normalfont\normalsize\bfseries}{\thesubsubsection}{1em}{}

% Espacements des listes
\setlist[itemize]{topsep=3pt, itemsep=2pt, parsep=0pt, leftmargin=1.5em}
\setlist[enumerate]{topsep=3pt, itemsep=2pt, parsep=0pt, leftmargin=1.5em}

\begin{document}

% --- PAGE DE GARDE ---
\begin{titlepage}
    \centering
    \vspace*{\fill}
    \rule{\linewidth}{1.2pt} \\[0.4cm]
    {\huge \bfseries \color{enstablu} Spécifications fonctionnelles \\[0.3cm] Cahier des charges du système informatique unifié des Foy'z et Bar \par}
    \rule{\linewidth}{1.2pt} \\[1.5cm]
    
    \vspace*{\fill}
    {\large Version 1.0 --- \today \par}
\end{titlepage}

\clearpage
\tableofcontents
\clearpage


\section{Introduction et contexte}
Le présent document a pour vocation de décrire de manière exhaustive les besoins fonctionnels du nouveau système informatique unifié destiné à la gestion des Foy’z/Bar. 

Dans le cadre de la fusion entre l'ENSTA Bretagne et l'ENSTA Paris, ce projet s'inscrit dans la démarche d'uniformisation des associations et de leurs plateformes. L'objectif est de remplacer les solutions locales d'administration des Foy'z/Bar par une infrastructure centralisée unique, adaptée aux particularités des deux campus.

Le système envisagé repose sur une architecture de type client-serveur :
\begin{itemize}
    \item Interface utilisateur : un site internet intuitif et sécurisé, servant d'interface publique pour les étudiants et de point d'entrée pour les équipes Foy'z/Bar.
    \item Base de données : une base de données SQL pour permettre notamment la traçabilité des soldes et des transactions et obtenir un historique des connexions.
\end{itemize}

Les soldes et sommes d'argent traitées sur le site ne sont pas directement reliés au compte de la trésorerie de l'organisme. La monnaie du système est donc une monnaie "fictive", censée faire le lien entre le suivi des consommations, le solde des usagers et la trésorerie réelle. Il va sans dire qu'un tel système repose sur l'honnêteté et la rigueur des équipes.

Il existe une distinction entre les prix standards et les prix équipe (pour les membres de l'équipe). En effet, une remise peut être appliquée sur certains articles pour les membres de l'équipe.

Chaque étudiant possède deux porte-monnaie virtuels distincts, un pour Brest et un pour Paris. Il n'est pas possible de transférer de l'argent du porte-monnaie d'un campus à l'autre et une consommation sur un campus ne peut être réglée à l'aide du porte-monnaie de l'autre campus.

Les statuts \textsc{"blacklist"} et \textsc{"blacklist alcool"} sont mutualisés entre les deux campus. Un individu ayant été membre dans l'équipe perd tous ses accès si il est désigné comme \textsc{blacklist}. Si il est désigné \textsc{blacklist alcool}, il conserve ses accès mais ne peut pas se retirer le statut \textsc{blacklist alcool}.

Le système web devra être protégé contre les potentielles tentatives d'intrusion extérieures.

\section{Interface publique}
La page d'accueil constitue l'interface publique du site internet. Elle doit être accessible à tout utilisateur sans authentification (étudiants, visiteurs extérieurs...). Elle regroupe les onglets suivants :

\begin{itemize}
    \item \textbf{Page d'accueil :} premier onglet visible lors de la visite du site. On peut y voir les événements prévus dans chaque campus au Foy'z/Bar et leur affiche. On peut aussi voir des notes publiques laissées par les équipes à l'attention des visiteurs du Foy'z/Bar et du site internet.
    \item \textbf{Accès aux prix standards :} consultation publique du catalogue des articles disponibles au bar avec affichage du prix standard en fonction du campus.
    \item \textbf{Règlement intérieur :} onglet dédié à la consultation des règlements intérieurs (Brest et Paris) du Foy'z/Bar. La dernière version des règlements doit pouvoir être consultée facilement.
    \item \textbf{Liens utiles :} cet onglet permet de donner plus de visibilité à des des liens d'utilité (par exemple les liens vers les plateformes de signalements VSS). Ces liens sont configurable dans le sous-onglet Module de développement de l'onglet administrateur.
    \item \textbf{Trombinoscopes :} présentation visuelle et nominative des deux équipes de mandat.
    \item \textbf{Page d'identification :} authentification sécurisée (Identifiant / Mot de passe) restreinte aux membres actifs et anciens membres de l'équipe pour accéder à l'interface équipe.
\end{itemize}

% --- SECTION 3 ---
\section{Interface équipe}
L'accès à l'interface équipe est exclusivement réservé aux membres et anciens membres des équipes Foy'z et Bar. Pour que les encaissements se fassent dans le bon porte-monnaie et que la mise-à-jour des bilans statistiques et de trésorerie soit exacts, le lieu de la connexions doit pouvoir être précisé entre Brest et Paris. Cette interface centralise la gestion quotidienne du Foy'z/Bar via les onglets suivants :

\begin{itemize}
    \item \textbf{Paiement :} interface de saisie rapide permettant d'enregistrer la commande d'un étudiant. Le membre de l'équipe saisit le nom de l'étudiant et sélectionne une liste d'articles. Des boutons \textbf{+} et \textbf{-} permettent d'ajuster la quantité facilement. À la validation, le solde du compte de l'étudiant est automatiquement débité. 
    
    Une commande doit pouvoir être partagée par plusieurs étudiants. Une option permet d'ajouter des contributeurs supplémentaires et le montant total de la commande est divisé équitablement entre tous les contributeurs. Si au moins un des contributeurs a le statut \textsc{"blacklist"}, la transaction ne peut pas aboutir et le système le notifie au membre de l'équipe. Si la commande contient ne serait-ce qu'un article alcoolisé et que ne serait-ce qu'un contributeur a le statut \textsc{"blacklist alcool"}, la transaction n'aboutit pas non plus et le système le notifie au membre de l'équipe. Si le solde d'un des contributeurs est plus faible que sa part à payer sur la commande, le système demande confirmation à l'aide du mot de passe administrateur pour valider la commande et ainsi placer l'étudiant correspondant dans le négatif. Si le montant a débiter creuse le découvert d'un ou de plusieurs contributeurs au delà de la limite autorisée, le système refuse la transaction. 
    
    Enfin, cet onglet doit présenter une option "payer directement" qui permet à un visiteur ne possédant pas de compte de régler directement sa commande. Ces transactions seront détaillées séparément dans les onglets Statistiques et Trésorerie.
    \item \textbf{Gestion des consignes :} intégrée à l'onglet de paiement précédent, cette option permet de suivre et de modifier le nombre de verres consignés non-rendus. Le mécanisme de consigne est activable sur demande et n'est pas systématique. La somme (réglable dans les paramètres) est prélevée sur le solde de l'étudiant à l'emprunt, puis créditée instantanément lors de la restitution du verre.
    \item \textbf{Rechargement :} permet de créditer le compte d'un étudiant. Le membre de l'équipe renseigne le nom, le montant à ajouter, le campus et le moyen de paiement utilisé parmi les quatre options supportées : Carte Bancaire, Lydia, Espèces ou HelloAsso.
    \item \textbf{Retrait :} utilisé lorsqu'un étudiant souhaite récupérer de l'argent liquide depuis son compte Foy'z/Bar. Le membre de l'équipe saisit le nom et le montant et sélectionne le le campus. Le système bloque la transaction si le montant demandé est supérieur au solde disponible du compte.
    \item \textbf{Transfert :} permet des virements entre deux étudiants et entre leur porte-monnaie de même campus. Le membre de l'équipe indique le compte donneur, le compte receveur et le montant. Une vérification s'assure que le solde du donneur est suffisant avant de procéder à la mise à jour du solde des deux comptes.
    \item \textbf{Historique des transactions :} journal global des transactions affichant les transactions par ordre chronologique inversé (la dernière en premier). Chaque entrée détaille : la date et l'heure, l'identité de l'étudiant, le type d'opération (Achat, Retrait, Transfert, Rechargement) et la somme concernée. Un système de filtres et une barre de recherche permettent d'affiner les recherches. Il doit être possible d'annuler une ou des transactions sur saisie du mot de passe administrateur. Suite à une annulation de transaction, les soldes concernés sont mis à jour en conséquence.
    \item \textbf{Statistiques de consommation :} tableau de bord présentant les volumes de vente et les tendances de consommation. Les données doivent pouvoir être filtrées par période temporelle, catégories d'articles, promotion d'étudiants, campus, statut Foy'z/Bar ou non etc.
    \item \textbf{Suivi trésorerie :} cet onglet permet de consulter les bilans financiers mensuels des campus. Ces bilans doivent renseigner les rechargements (somme et détail du type de paiement), les ventes (valeur et détail selon le type d'article) ainsi que les recettes des évènements. Un graphe permet de visualiser facilement ces données sur douze mois glissants.
    \item \textbf{Systèmes de notes (Post-it) :} 
    \begin{itemize}
        \item un espace collaboratif privé servant de tableau d'affichage virtuel pour les équipes. Tout membre peut publier un message informatif, modifier une note existante ou la supprimer. Afin d'éviter l'encombrement visuel, seules les notes les plus récentes sont maintenues à l'écran, avec un nombre maximal de notes à paramétrer.
        \item un espace collaboratif public permettant aux équipes de communiquer des informations à l'intention des visiteurs sur la page d'accueil du site. Seuls les membres peuvent publier des messages.
    \end{itemize}
\end{itemize}

% --- SECTION 4 ---
\section{Sous-onglet d'administration}
L'onglet \textbf{Administrateur} est l'interface de contrôle du système. Il englobe plusieurs sous-onglets techniques destinés à la maintenance du site :

\begin{itemize}
    \item \textbf{Gestion des comptes :} moteur de recherche global sur les profils étudiants. Il donne accès à la totalité des informations de chaque utilisateur : nom, solde, ainsi que ses indicateurs de restriction : Statut \textsc{blacklist} et \textsc{blacklist alcool}. Les profils sont modifiables à tout moment par les membres de l'équipe.
    \item \textbf{Equipe :} cet onglet permet de gérer les statuts \textsc{équipe} qui indique si un compte possède les accès équipe. Il est aussi possible de préciser si le membre de l'équipe est de mandat ou est un ancien membre. Les anciens membres conservent les accès aux onglets de l'interface équipe mais perdent l'accès aux onglets administrateurs.
    \item \textbf{Gestion des articles :} interface dédiée au catalogue des articles. Chaque fiche article comprend : le nom, le prix standard, le prix Foy'z/barreux (pour chaque campus), le volume (pour les boissons) et le type d'article (parmi Bière, Vin, Cidre, Snacks, Saucisson, Évènement/Soirées). Les articles de tireuse sont mis à jour dans l'onglet suivant et ne sont donc pas modifiables dans cet onglet. De même, les articles temporaires liés à des évènements sont uniquement modifiables dans l'onglet Gestion des évènement.
    \item \textbf{Bières pression et tireuses :} onglet permettant le suivi des fûts enregistrés (Nom, degré d'alcool, volume total, prix au demi, à la pinte et au pot pour les prix standards et équipe et selon les campus). Cet onglet permet également d'assigner un fût à une tireuse. Cette association met automatiquement à jour le catalogue des articles tireuse en respectant le format : \texttt{"[Demi/Pinte/Pot] de tireuse [Numéro] ["Nom du Fût"]"}
    \item \textbf{Gestion des événements :} permet de planifier des soirées spéciales en définissant un nom, des heures de début et de fin et un catalogue d'articles temporaires associés. Ce module génère une interface simplifiée d'encaissement destinée aux associations organisatrices (comme le BDE), leur permettant de gérer leurs ventes durant leurs évènements sans avoir accès à l'ensemble du système d'administration du Foy'z/Bar. Cette page est accessible via une passerelle à sens unique disponible dans cet onglet. Elle se résume à l'onglet paiement et permet d'encaisser les articles temporaires de l'évènement ainsi que le catalogue standard. Une nouvelle authentification est demandée pour revenir à l'interface équipe.
    \item \textbf{Registre des connexions :} traçabilité enregistrant les connexions des membres de l'équipe (Identifiant du compte, date et heure). L'interface propose une recherche par filtre par date et par nom. Une suppression  de ces logs est prévue au bout d'une durée configurable.
    \item \textbf{Module développement :} onglet technique qui regroupe la configuration des variables globales (paramètres) du site et fournit les accès directs (liens sécurisés) vers l'hébergeur, le serveur de base de données et le dépôt du code source.
    Parmi les paramètres configurables doivent au moins figurer les paramètres suivants:
    \begin{itemize}
        \item Montant limite d'autorisation de découvert.
        \item Valeur de la consigne.
        \item Fichiers PDF des Règlements intérieurs des deux campus pour l'onglet correspondant dans l'interface publique.
        \item Édition de la page d'accueil de l'interface publique décrite plus haut.
        \item Possibilité de modifier le thème de couleur et le logo du site indépendamment pour les deux campus.
        \item Édition des trombinoscopes.
        \item Configuration des durées maximales d'accès aux données des transactions depuis le site. Les données sont toujours conservées dans la base de données mais plus accessibles depuis le site après le délai fixé.
        \item Configuration du délai de déconnexion pour inactivité des sessions.
        \item Configuration du mot de passe administrateur.
        \item Configuration du nombre maximal de post-it affichés dans les notes.
        \item Configuration des liens utiles de l'interface publique.
    \end{itemize}
\end{itemize}

\end{document}